\documentclass[10pt,a4paper]{amsart}
\usepackage[margin=19mm]{geometry}
\usepackage{amsmath,amssymb,mathtools}
\usepackage[scr=rsfs,cal=euler]{mathalfa}
\usepackage{xfrac}
\usepackage[unicode,hidelinks,bookmarks=false]{hyperref}
\newcommand{\ie}{i.e.\ }
\newcommand{\cf}{cf.\ }
\newcommand{\resp}{respectively\ }
\newcommand{\Subsection}{\subsection}
\providecommand{\nobreakdash}{\nobreak\hbox{-}}
\setlength{\emergencystretch}{2em}
\multlinegap0pt
\usepackage{cancel}
\usepackage{tikz}
\usepackage{amssymb}
\usepackage{xcolor}
\usepackage[normalem]{ulem}
\newtheorem{lem}{Lemma}
\newcommand{\BB}{\mathfrak{B}}
\newcommand{\Det}{\operatorname{Det}}
\newcommand{\ad}{\mathrm{ad}}
\newcommand{\bt}{\bullet}
\newcommand{\im}{\operatorname{im}}
\newcommand{\la}{\leftarrow}
\newcommand{\mc}{\mathcal}
\newcommand{\mr}{\mathrm}
\newcommand{\sdet}{\operatorname{Sdet}}
\title{Very important paper: On the dependence of the perturbative Chern-Simons partition function on the reference flat connection in the BV formalism}
\author{Pavel Mnev, Konstantin Wernli}
\date{Source: arXiv:2510.18653v3 (CC BY 4.0)}
\begin{document}
\maketitle

\noindent\emph{Source and licence.} Very important paper: On the dependence of the perturbative Chern-Simons partition function on the reference flat connection in the BV formalism. Version arXiv:2510.18653v3 (CC BY 4.0), distributed by arXiv under the Creative Commons Attribution 4.0 International licence, http://creativecommons.org/licenses/by/4.0/, as recorded in the arXiv licence metadata for this exact version and shown on the abstract page https://arxiv.org/abs/2510.18653v3. Full licence notice: \texttt{licenses/arxiv-2510.18653-CC-BY-4.0.txt}.\par\smallskip

\noindent\textit{Editorial scope.} Counted unit: the article's lem A' (source0.tex lines 2700--2706). The article's complete original proof of it is delivered with it (source0.tex lines 2755--2814, one \texttt{proof} environment of 60 lines, which the article places away from the statement). No mathematics is written, repaired or re-derived: the statement and the proof are the article's own lines, in the article's own notation, and no retained line is substituted or re-flowed.  Retained uncounted as the source-local context the proof needs: lines 2200--2204; lines 2631--2635; lines 2650--2654; lines 2714--2719.  Nothing was omitted from any retained span, so the package carries no omission record.  No retained line contains a citation, so the wrapper carries no bibliography and no bibliography entry.  No retained line contains a figure environment, an image-inclusion command or drawing code, so no image is delivered and the package declares no asset dependency.  Editorial formatting only, in addition: an AMS \texttt{amsart} preamble that restates the article's own declarations and macros used by the retained lines, namely the environments \texttt{lem} on separate counters, together with the standard packages those lines need.  The retained environments print in this document as Lemma 1 in order of appearance, whereas the article prints the same items with the numbering of its own full document.  The complete native source of the article as distributed in the arXiv e-print for this exact version is bundled unchanged as \texttt{sources/187514/source0.tex}.
\begin{equation}\label{I def}
    I_{A,A'} = e^{\frac{i\pi}{4} \psi_{A'}}%\det(q)^{-1/2}
    \det(\BB_{A\la A';A'})^{1/2}
    \tau_{A'}^{1/2}\sdet^{\frac{1}{2}}_{\im d_{A'}^*}\left(1 + K_{A'}\operatorname{ad}_\beta\right) \in \Det^{\frac12} H^\bullet_{A},
\end{equation}

\begin{equation}\label{Lambda, II, PP}
    \Lambda = K_{A,A'} \mr{ad}^*_{\delta A'} G_{A,A'}, \quad
    \mathbb{I} = G_{A,A'} \mr{ad}^*_{\delta A'}, \quad
    \mathbb{P} = \mr{ad}^*_{\delta A'} G_{A,A'}.
\end{equation}

\begin{multline}\label{Xi}
    \Xi=p \Big( \mathbb{P}\,\ad_\beta (1-K\ad_\beta)^{-1}-(1-\ad_\beta K)^{-1} \ad_\beta\, \mathbb{I} +\\
    +  (1-\ad_\beta K)^{-1}\ad_\beta\, \Lambda\, \ad_\beta (1-K\ad_\beta)^{-1} \Big) i
%   \Xi=p_{A',A'} \Big( \mathbb{P}\ad_\beta -\ad_\beta \mathbb{I}+ \ad_\beta \Lambda \ad_\beta \Big) i_{A',A'}
\end{multline}

\begin{lem}\label{lemma: dependence of torsion on A}
Given a path of flat connections $A_t$, one has the following formula for the infinitesimal change of the Ray-Singer torsion $\tau_{A_t}$:
    \begin{equation}\label{tau dependence on A}
        \left.\frac{d}{dt}\right|_{t=0}\det(\BB_{A_0\la A_t})\tau_{A_t}=\mr{Str}_{\Omega^\bt}(K_{A_0}\ad_{\dot{A}_0})\, \tau_{A_0}.
    \end{equation}
\end{lem}

\begin{lem}\label{lemma: dependence of I on A'}
   The variation of $I_{A,A'}$ w.r.t. $A'$ is given by
   \begin{equation}
       \delta_{A'} I_{A,A'}= \Xi\, I_{A,A'}.
   \end{equation}
%   with $\Xi$ as in (\ref{Xi}).
\end{lem}

\begin{proof}[Proof of Lemma \ref{lemma: dependence of I on A'}]
    From the definition (\ref{I def}) of $I_{A,A'}$ and from Lemma \ref{lemma: dependence of torsion on A} we find
\begin{multline}\label{lem I eq1}
    2I_{A,A'}^{-1}  \delta_{A'} I_{A,A'} =\underbrace{\mr{Str}_{\Omega^\bt}\cancel{K_{A'}\ad_{\delta A'}}}_{\mr{from}\; (\ref{tau dependence on A})}
    +
    \mr{Str}_{\Omega^\bt} (1+K_{A'}\ad_\beta)^{-1}\Big(\underbrace{-\cancel{K_{A'}\ad_{\delta A'}}}_{K_{A'}\delta_{A'}\ad_\beta}-\\
    \underbrace{-\cancel{K_{A'}\ad_{\delta A'}K_{A'} \ad_\beta} + [d_{A'},K_{A'}\ad^*_{\delta A'}G_{A'}] \ad_\beta +P_{A'}\ad^*_{\delta A'}G_{A'} \ad_\beta+G_{A'}\ad^*_{\delta A'} P_{A'} \ad_\beta}_{(\delta_{A'}K_{A'})\ad_\beta}\Big)
 %   \\= \mr{Str}_{Omega^\bt} -\ad_\beta (1+K_{A'}\ad_\beta)^{-1}[d_{A'},]
\end{multline}
Recall that $\beta=A-A'$. Denote $\lambda=K_{A'}\ad^*_{\delta A'}G_{A'}$.
Focusing on the first non-cancelling term above, we  find
\begin{multline}\label{lem I eq2}
\mr{Str} (-\ad_\beta (1+K\ad_\beta)^{-1} [d,\lambda])
=\mr{Str}(-[d,\ad_\beta (1+K\ad_\beta)^{-1}]\lambda) \\=
\mr{Str} \Big(\underbrace{-[d,\ad_\beta]}_{(\ad_\beta)^2} (1+K\ad_\beta)^{-1}-\\
-\ad_\beta (1+K\ad_\beta)^{-1} \underbrace{[d,K\ad_\beta]}_{(1-P)\ad_\beta+K (\ad_\beta)^2 =
(1+K\ad_\beta)\ad_\beta-P\ad_\beta} (1+K\ad_\beta)^{-1} \Big)\lambda\\
=\mr{Str}\, \ad_\beta (1+K\ad_\beta)^{-1}   P \ad_\beta (1+K\ad_\beta)^{-1} \lambda
\end{multline}
Here we were suppressing the subscripts $A'$ in $d_{A'},K_{A'},P_{A'}$. Also we used that, due to flatness of $A$ and $A'$, we have $[d,\ad_\beta]=-(\ad_\beta)^2$.
Next, we notice that $(1+K\ad_\beta)^{-1}K=K_{A,A'}$ and that
\begin{align*}
    G_{A,A'}=\sum_{k\geq 0}(-1)^k (G[\ad_\beta,d^*]_+)^k G &= &G(1+\ad_\beta K)^{-1}+d^* (\cdots)\\
    &=& (1+K\ad_\beta)^{-1}G+(\cdots)d^*
\end{align*}
This implies
\begin{equation}
    (\ref{lem I eq2})=-\mr{Str}\, P_{A'} \ad_\beta K_{A,A'} \ad^*_{\delta A'}G_{A,A'}\ad_\beta
\end{equation}
Thus, returning to the computation (\ref{lem I eq1}), we have
\begin{multline}\label{lem I eq3}
 2I_{A,A'}^{-1}  \delta_{A'}I_{A,A'}=\\
 =
 \mr{Str}_{\Omega^\bt} P_{A'}\Big(-\ad_\beta K_{A,A'} \ad^*_{\delta A'}G_{A,A'}\ad_\beta-\ad_\beta G_{A,A'} \ad^*_{\delta A'}+\ad^*_{\delta A'}G_{A,A'}\ad_\beta
 \Big)
 \\ =
 \mr{Str}_{\Omega^\bt}P_{A'}(-\ad_\beta \Lambda \ad_\beta - \ad_\beta \mathbb{I} + \mathbb{P}\ad_\beta)=2\Xi
\end{multline}
with $\Lambda,\mathbb{I},\mathbb{P}$ as in (\ref{Lambda, II, PP}) and $\Xi$ as defined in (\ref{Xi}).
% \begin{comment}
% Next, remark that by homological perturbation theory formulae for the deformation of $(i,p,K)$ triples, we have
% \begin{equation*}
%     i_{A',A'}=\underbrace{(1-K_{A,A'}\ad_\beta)^{-1}}_{\mc{L}}i_{A,A'}\BB_{A\la A'},\;\; p_{A',A'}=\BB_{A'\la A}\,p_{A,A'}\underbrace{(1-\ad_\beta K_{A,A'})^{-1}}_{\mc{R}}
% \end{equation*}
% Therefore,
% \begin{equation*}
%     P_{A'}=i_{A',A'}p_{A',A'}=
%     %(1-K_{A,A'}\ad_\beta)^{-1}\underbrace{i_{A,A'}p_{A,A'}}_{P_{A,A'}} (1-\ad_\beta K_{A,A'})^{-1}
%     \mc{L}\ i_{A,A'}p_{A,A'}\ \mc{R}= \mc{L} P_{A,A'}\mc{R}
% \end{equation*}
% Using this, the projection onto $(A',A')$-harmonic forms in (\ref{lem I eq3}) can be traded for the projection onto $(A,A')$-harmonic forms:
% \begin{multline*}
%  2I_{A,A'}^{-1}  \delta_{A'}I_{A,A'}=\mr{Str}_{\Omega^\bt} P_{A,A'} \mc{R} (\ad_\beta \Lambda \ad_\beta
%  -\ad_\beta\mathbb{I} (1-K_{A,A'}\ad_\beta)+(1-\ad_\beta K_{A,A'})\mathbb{P}\ad_\beta) \mc{L}   \\
%  =\mr{Str}_{\Omega^\bt}P_{A,A'} (\mc{R}\ad_\beta \Lambda \ad_\beta\mc{L}-\mc{R} \ad_\beta\mathbb{I}+ \mathbb{P}\ad_\beta \mc{L})=
%  \mr{Str}_{H^\bt_A} \Xi
% \end{multline*}
% with $\Xi$ as in (\ref{Xi}).
% \end{comment}
\end{proof}

\end{document}
